\scanpage{302}
\mspage{1}
\margin{1975\quad\unsure{à}\quad1976}
\[
\boldsymbol\Delta=\text{cat. des simplexes }\Delta^n\quad(n\geq0)
\]
\margin{$\Delta^n=[0,n]\subset\mathbf N$\quad avec rel. d'ordre \unsure{total|totale}}
\[
\widehat{\boldsymbol\Delta}=\mathrm{Ss}=\underline{\operatorname{Hom}}(\boldsymbol\Delta^o,\mathrm{Ens})=\mathrm{Ens}_*
\]
Plus généralement, si $A$ est une catégorie, on pose
\[
\begin{aligned}
A_*&=\underline{\operatorname{Hom}}(\boldsymbol\Delta^o,\mathrm{Ens})\\
A^*&=\underline{\operatorname{Hom}}(\boldsymbol\Delta,\mathrm{Ens})
\end{aligned}
\]
donc
\[ (A^o)^*=(A_*)^o\qquad A^*\simeq((A^o)_*)^o \]
où l'exposant $o$ désigne le \inserted{\unsure{passage à la}} cat. opposée.

Les objets de $A_*$ (resp. $A^*$) sont considérés comme des structures algébriques de type simple (sur une infinité d'objets de base) dans $A$, les ``objets $\tfrac12$ simpliciaux'' resp. ``$\tfrac12$ cosimpliciaux''. On écrira $\mathrm{Ens}_*$ quand on a en vue cet aspect, et $\widehat{\boldsymbol\Delta}$ \inserted{ou $\mathrm{Ss}$} quand on a plutôt le \unsure{p\textsuperscript{t}} de vue ``objet d'un topos'' ou ``faisceau sur un topos'' $\unsure{\flat}$. Tous les développements qui suivront ont pour objet d'étudier ces objets (l'équivalent combinatoire des espaces topologiques) et leurs ``types d'homotopie'', \unsure{i.e.} des invariants qu'on peut leur associer, qui en dépendent de façon covariante ou contravariante. Nous avons ici en vue une étude systématique, dans l'esprit de l'algèbre universelle, de ces invariants, (qui sont donc des foncteurs sur $\widehat{\boldsymbol\Delta}\simeq\mathrm{Ens}_*$) \deleted{et des} \inserted{de leurs structures et des} opérations qui peuvent être définies sur elles, permettant d'en déduire certains complexes \unsure{à} partir \unsure{d'autres} plus simples

\margin{$\unsure{\flat}$ La notation $\widehat{\boldsymbol\Delta}$ peut être utilisée de préf. quand on veut garder à l'esprit la nature \uline{algébrique} très particulière de ces \unsure{types}, la notation $\mathrm{Ss}$ quand on veut l'oublier au profit de la \uline{signification} topologique.}
\margin{303}
